\documentclass{article}
\usepackage{amsmath, amssymb, amsthm}
\usepackage{enumitem}
\usepackage[utf8]{inputenc}
\usepackage[spanish]{babel}
\usepackage{mathrsfs}
\usepackage{mathtools}
\usepackage{geometry}
\geometry{a4paper, margin=1in}

\theoremstyle{plain}
\newtheorem{Teorema}{Teorema}[section]
\newtheorem{Proposicion}[Teorema]{Proposici\'on}
\newtheorem{Lema}[Teorema]{Lema}
\newtheorem{Corolario}[Teorema]{Corolario}

\theoremstyle{definition}
\newtheorem{Definicion}[Teorema]{Definici\'on}
\newtheorem{Ejemplo}[Teorema]{Ejemplo}
\newtheorem{Observacion}[Teorema]{Observaci\'on}
\newtheorem{Ejercicio}{Ejercicio}[section]

\newcommand{\N}{\mathbb{N}}
\newcommand{\Z}{\mathbb{Z}}
\newcommand{\R}{\mathbb{R}}
\newcommand{\C}{\mathbb{C}}
\newcommand{\E}{\mathbb{E}}
\newcommand{\Prob}{\mathbb{P}}
\newcommand{\A}{\mathbf{A}}
\newcommand{\I}{\mathbf{I}}
\newcommand{\p}{\mathbf{p}}
\newcommand{\PP}{\mathbf{P}}
\newcommand{\ind}{\mathbf{1}}  % Fixed: was \1 which is invalid

\title{Cadenas de Markov en Tiempo Continuo}
\author{Gal}
\date{}

\begin{document}

\maketitle

\section{Propiedades Fundamentales de las Distribuciones Exponenciales}

\begin{Definicion}
  Sean $T_1, \ldots, T_n$ variables aleatorias independientes con distribuci\'on exponencial de par\'ametros $b_1, \ldots, b_n$ respectivamente. Definimos
  \[ T = \min\{T_1, \ldots, T_n\}. \]
  Interpretamos $T_i$ como el momento en que suena el reloj $i$.
\end{Definicion}

\begin{Teorema}
  La variable aleatoria $T$ tiene distribuci\'on exponencial con par\'ametro $b_1 + \cdots + b_n$.
\end{Teorema}

\begin{proof}
  Calculamos la funci\'on de supervivencia:
  \[
  \Prob\{T \geq t\} = \Prob\{T_1 \geq t, \ldots, T_n \geq t\}.
  \]
  Por independencia:
  \[
  \Prob\{T \geq t\} = \prod_{i=1}^n \Prob\{T_i \geq t\}.
  \]
  Cada $T_i$ es exponencial de par\'ametro $b_i$, por lo que $\Prob\{T_i \geq t\} = e^{-b_i t}$. Sustituyendo:
  \[
  \Prob\{T \geq t\} = \prod_{i=1}^n e^{-b_i t} = e^{-(b_1 + \cdots + b_n)t}.
  \]
  Esta es precisamente la funci\'on de supervivencia de una variable exponencial con par\'ametro $\sum_{i=1}^n b_i$.
\end{proof}

\begin{Teorema}
  La probabilidad de que el reloj $i$ sea el primero en sonar es
  \[ \Prob\{T_i = T\} = \frac{b_i}{b_1 + \cdots + b_n}. \]
\end{Teorema}

\begin{proof}
  Condicionamos en el valor de $T_i$. Sea $f_{T_i}(t) = b_i e^{-b_i t}$ la densidad de $T_i$. Entonces:
  \[
  \Prob\{T_i = T\} = \int_0^\infty \Prob\{T_j > t \text{ para todo } j \neq i \mid T_i = t\} \, f_{T_i}(t) \, dt.
  \]
  Por independencia, condicionar en $T_i = t$ no afecta a los otros $T_j$:
  \[
  \Prob\{T_j > t \text{ para todo } j \neq i \mid T_i = t\} = \prod_{j \neq i} \Prob\{T_j > t\} = \prod_{j \neq i} e^{-b_j t} = e^{-(b_1 + \cdots + b_{i-1} + b_{i+1} + \cdots + b_n)t}.
  \]
  Por tanto:
  \[
  \Prob\{T_i = T\} = \int_0^\infty e^{-(b_1 + \cdots + b_{i-1} + b_{i+1} + \cdots + b_n)t} \cdot b_i e^{-b_i t} \, dt = b_i \int_0^\infty e^{-(b_1 + \cdots + b_n)t} \, dt.
  \]
  La integral es $\int_0^\infty e^{-Bt} \, dt = 1/B$ con $B = b_1 + \cdots + b_n$. Por tanto:
  \[
  \Prob\{T_i = T\} = \frac{b_i}{b_1 + \cdots + b_n}.
  \]
\end{proof}

\section{Cadenas de Markov en Tiempo Continuo}

\begin{Definicion}
  Sea $S$ un espacio de estados finito. Un proceso en tiempo continuo $\{X_t\}_{t \geq 0}$ con valores en $S$ es una \textbf{cadena de Markov en tiempo continuo} si satisface:
  \begin{enumerate}
    \item \textbf{Propiedad de Markov:} Para todo $t > s$,
    \[ \Prob\{X_t = y \mid X_r, 0 \leq r \leq s\} = \Prob\{X_t = y \mid X_s\}. \]
    \item \textbf{Homogeneidad en el tiempo:} Para todo $t > s$,
    \[ \Prob\{X_t = y \mid X_s = x\} = \Prob\{X_{t-s} = y \mid X_0 = x\}. \]
  \end{enumerate}
\end{Definicion}

\begin{Definicion}
  Para cada par de estados $x, y \in S$ con $x \neq y$, asignamos un n\'umero no negativo $\alpha(x,y)$ llamado \textbf{tasa de transici\'on} de $x$ a $y$. Definimos la \textbf{tasa total de salida} del estado $x$ como
  \[ \alpha(x) = \sum_{y \neq x} \alpha(x,y). \]
\end{Definicion}

\begin{Definicion}
  La cadena de Markov en tiempo continuo con tasas $\alpha$ satisface las siguientes probabilidades infinitesimales:
  \begin{align}
    \Prob\{X_{t+\Delta t} = x \mid X_t = x\} &= 1 - \alpha(x)\Delta t + o(\Delta t), \label{eq:stay}\\
    \Prob\{X_{t+\Delta t} = y \mid X_t = x\} &= \alpha(x,y)\Delta t + o(\Delta t), \quad y \neq x. \label{eq:jump}
  \end{align}
  Aqu\'i $o(\Delta t)$ denota una funci\'on tal que $\lim_{\Delta t \to 0} o(\Delta t)/\Delta t = 0$.
\end{Definicion}

\subsection{Ecuaciones Diferenciales y el Generador Infinitesimal}

\begin{Definicion}
  Sea $p_x(t) = \Prob\{X_t = x\}$.
\end{Definicion}

\begin{Proposicion}
  Las probabilidades $p_x(t)$ satisfacen el sistema de ecuaciones diferenciales:
  \[ p_x'(t) = -\alpha(x)p_x(t) + \sum_{y \neq x} \alpha(y,x)p_y(t). \]
\end{Proposicion}

\begin{proof}
  Consideramos $p_x(t+\Delta t)$ y condicionamos en el estado en tiempo $t$:
  \[
  p_x(t+\Delta t) = \sum_{y \in S} \Prob\{X_t = y\} \Prob\{X_{t+\Delta t} = x \mid X_t = y\}.
  \]
  Usando las probabilidades infinitesimales \eqref{eq:stay} y \eqref{eq:jump}:
  \begin{align*}
    p_x(t+\Delta t) &= p_x(t)(1 - \alpha(x)\Delta t + o(\Delta t)) + \sum_{y \neq x} p_y(t)(\alpha(y,x)\Delta t + o(\Delta t)) \\
    &= p_x(t) - \alpha(x)p_x(t)\Delta t + \sum_{y \neq x} \alpha(y,x)p_y(t)\Delta t + p_x(t)o(\Delta t) + \sum_{y \neq x} p_y(t)o(\Delta t).
  \end{align*}
  Reordenamos:
  \[
  \frac{p_x(t+\Delta t) - p_x(t)}{\Delta t} = -\alpha(x)p_x(t) + \sum_{y \neq x} \alpha(y,x)p_y(t) + \frac{o(\Delta t)}{\Delta t}\left(p_x(t) + \sum_{y \neq x} p_y(t)\right).
  \]
  El \'ultimo t\'ermino tiende a 0 cuando $\Delta t \to 0$ porque $o(\Delta t)/\Delta t \to 0$ y la suma de probabilidades es finita. Por tanto, tomando el l\'imite:
  \[
  p_x'(t) = \lim_{\Delta t \to 0} \frac{p_x(t+\Delta t) - p_x(t)}{\Delta t} = -\alpha(x)p_x(t) + \sum_{y \neq x} \alpha(y,x)p_y(t).
  \]
\end{proof}

\begin{Definicion}
  La \textbf{matriz generadora infinitesimal} $\A$ de la cadena tiene entradas:
  \[ \A(x,y) = \begin{cases}
    \alpha(x,y) & \text{si } x \neq y,\\
    -\alpha(x) & \text{si } x = y.
  \end{cases} \]
\end{Definicion}

\begin{Observacion}
  Las filas de $\A$ suman cero porque:
  \[ \sum_{y \in S} \A(x,y) = -\alpha(x) + \sum_{y \neq x} \alpha(x,y) = -\alpha(x) + \alpha(x) = 0. \]
  Adem\'as, las entradas no diagonales son no negativas y las diagonales son no positivas.
\end{Observacion}

\begin{Proposicion}
  Sea $\p(t)$ el vector fila de probabilidades, es decir, $\p(t) = (p_x(t))_{x \in S}$. Entonces
  \[ \p'(t) = \p(t)\A. \]
\end{Proposicion}

\begin{proof}
  La componente $x$ de $\p(t)\A$ es:
  \[ (\p(t)\A)_x = \sum_{y \in S} p_y(t) \A(y,x) = p_x(t)\A(x,x) + \sum_{y \neq x} p_y(t) \A(y,x). \]
  Sustituyendo $\A(x,x) = -\alpha(x)$ y $\A(y,x) = \alpha(y,x)$ para $y \neq x$:
  \[ (\p(t)\A)_x = -p_x(t)\alpha(x) + \sum_{y \neq x} p_y(t)\alpha(y,x) = p_x'(t). \]
\end{proof}

\begin{Teorema}
  La soluci\'on del sistema $\p'(t) = \p(t)\A$ con condici\'on inicial $\p(0)$ es
  \[ \p(t) = \p(0)e^{t\A}, \]
  donde $e^{t\A} = \sum_{n=0}^\infty \frac{(t\A)^n}{n!}$ es la exponencial de matrices.
\end{Teorema}

\begin{proof}
  Verificamos que $\p(t) = \p(0)e^{t\A}$ satisface la ecuaci\'on diferencial:
  \[ \frac{d}{dt} \p(t) = \p(0) \frac{d}{dt} e^{t\A} = \p(0) \A e^{t\A} = (\p(0)e^{t\A}) \A = \p(t) \A. \]
  La condici\'on inicial se cumple porque $e^{0\A} = \I$.
\end{proof}

\begin{Definicion}
  Sea $\PP_t$ la matriz de transici\'on con entradas $p_t(x,y) = \Prob\{X_t = y \mid X_0 = x\}$.
\end{Definicion}

\begin{Proposicion}
  La matriz $\PP_t$ satisface la ecuaci\'on matricial
  \[ \frac{d}{dt} \PP_t = \PP_t \A, \quad \PP_0 = \I, \]
  y por tanto $\PP_t = e^{t\A}$.
\end{Proposicion}

\begin{proof}
  Fijamos $x$ y consideramos el vector $\p(t) = (\PP_t(x,\cdot))$, es decir, la fila $x$ de $\PP_t$. Por definici\'on, $\p(t)$ son las probabilidades condicionadas a $X_0 = x$, por lo que satisfacen $\p'(t) = \p(t)\A$. Esto es precisamente la fila $x$ de la ecuaci\'on $\frac{d}{dt}\PP_t = \PP_t\A$. La condici\'on inicial $\PP_0 = \I$ se cumple porque $p_0(x,y) = \delta_{xy}$.
\end{proof}

\begin{Ejemplo}[Cadena de dos estados]
  Consideremos una cadena con estados $\{0,1\}$, tasas $\alpha(0,1) = 1$, $\alpha(1,0) = 2$. El generador infinitesimal es
  \[ \A = \begin{pmatrix}
    -1 & 1 \\
    2 & -2
  \end{pmatrix}. \]
  Para calcular $e^{t\A}$, diagonalizamos la matriz. Los eigenvalores satisfacen:
  \[ \det(\A - \lambda \I) = \det\begin{pmatrix}
    -1-\lambda & 1 \\
    2 & -2-\lambda
  \end{pmatrix} = (-1-\lambda)(-2-\lambda) - 2 = \lambda^2 + 3\lambda = \lambda(\lambda+3) = 0. \]
  Por tanto, $\lambda_1 = 0$, $\lambda_2 = -3$.
  
  Para $\lambda = 0$, buscamos un eigenvector derecho $\A v = 0$:
  \[ \begin{pmatrix}
    -1 & 1 \\
    2 & -2
  \end{pmatrix} \begin{pmatrix} v_1 \\ v_2 \end{pmatrix} = \begin{pmatrix} 0 \\ 0 \end{pmatrix} \Rightarrow -v_1 + v_2 = 0 \Rightarrow v_2 = v_1. \]
  Tomamos $v = (1,1)^T$.
  
  Para $\lambda = -3$, resolvemos $(\A + 3\I)v = 0$:
  \[ \begin{pmatrix}
    2 & 1 \\
    2 & 1
  \end{pmatrix} \begin{pmatrix} v_1 \\ v_2 \end{pmatrix} = \begin{pmatrix} 0 \\ 0 \end{pmatrix} \Rightarrow 2v_1 + v_2 = 0 \Rightarrow v_2 = -2v_1. \]
  Tomamos $v = (1,-2)^T$.
  
  Construimos $Q = \begin{pmatrix} 1 & 1 \\ 1 & -2 \end{pmatrix}$ (eigenvectores como columnas). Entonces:
  \[ Q^{-1} = \frac{1}{\det(Q)} \begin{pmatrix} -2 & -1 \\ -1 & 1 \end{pmatrix} = -\frac{1}{3} \begin{pmatrix} -2 & -1 \\ -1 & 1 \end{pmatrix} = \frac{1}{3}\begin{pmatrix} 2 & 1 \\ 1 & -1 \end{pmatrix}. \]
  
  Verificamos que $Q^{-1} \A Q = D = \begin{pmatrix} 0 & 0 \\ 0 & -3 \end{pmatrix}$.
  
  Entonces:
  \[ e^{t\A} = Q e^{tD} Q^{-1} = Q \begin{pmatrix} 1 & 0 \\ 0 & e^{-3t} \end{pmatrix} Q^{-1}. \]
  
  Calculamos:
  \begin{align*}
    e^{t\A} &= \begin{pmatrix} 1 & 1 \\ 1 & -2 \end{pmatrix} \begin{pmatrix} 1 & 0 \\ 0 & e^{-3t} \end{pmatrix} \cdot \frac{1}{3}\begin{pmatrix} 2 & 1 \\ 1 & -1 \end{pmatrix} \\
    &= \frac{1}{3} \begin{pmatrix} 1 & e^{-3t} \\ 1 & -2e^{-3t} \end{pmatrix} \begin{pmatrix} 2 & 1 \\ 1 & -1 \end{pmatrix} \\
    &= \frac{1}{3} \begin{pmatrix}
      2 + e^{-3t} & 1 - e^{-3t} \\
      2 - 2e^{-3t} & 1 + 2e^{-3t}
    \end{pmatrix}.
  \end{align*}
  
  Podemos reescribir como:
  \[ \PP_t = \begin{pmatrix} \frac{2}{3} & \frac{1}{3} \\ \frac{2}{3} & \frac{1}{3} \end{pmatrix} + e^{-3t} \begin{pmatrix} \frac{1}{3} & -\frac{1}{3} \\ -\frac{2}{3} & \frac{2}{3} \end{pmatrix}. \]
  
  Observamos que $\lim_{t\to\infty} \PP_t = \begin{pmatrix} 2/3 & 1/3 \\ 2/3 & 1/3 \end{pmatrix}$. La distribuci\'on l\'imite es $\pi = (2/3, 1/3)$, que satisface $\pi \A = 0$:
  \[ \pi \A = \begin{pmatrix} \frac{2}{3} & \frac{1}{3} \end{pmatrix} \begin{pmatrix} -1 & 1 \\ 2 & -2 \end{pmatrix} = \begin{pmatrix} -\frac{2}{3} + \frac{2}{3} & \frac{2}{3} - \frac{2}{3} \end{pmatrix} = \begin{pmatrix} 0 & 0 \end{pmatrix}. \]
\end{Ejemplo}

\begin{Ejemplo}[Cadena de cuatro estados simétrica]
  Consideremos una cadena con estados $\{0,1,2,3\}$ y generador
  \[ \A = \begin{pmatrix}
    -3 & 1 & 1 & 1 \\
    2 & -4 & 1 & 1 \\
    2 & 1 & -4 & 1 \\
    1 & 1 & 1 & -3
  \end{pmatrix}. \]
  
  Para encontrar los eigenvalores, observamos que la matriz tiene estructura de bloque. Podemos encontrar un eigenvector $(1,1,1,1)$ con eigenvalor 0 porque la suma de filas es 0. Por simetr\'ia, buscamos vectores de la forma $(1,0,-1/2,-1/2)$, $(0,0,-1/2,1/2)$ y $(-1/3,1,-1/3,1/3)$. Estos son linealmente independientes y se puede verificar que son eigenvectores con eigenvalores $-1$, $-3$ y $-4$ respectivamente.
  
  La matriz $Q$ de eigenvectores es:
  \[ Q = \begin{pmatrix}
    1 & 1 & 0 & -1/3 \\
    1 & 0 & 0 & 1 \\
    1 & -1/2 & -1/2 & -1/3 \\
    1 & -1/2 & 1/2 & 1/3
  \end{pmatrix}. \]
  
  Entonces $D = Q^{-1}\A Q = \begin{pmatrix}
    0 & 0 & 0 & 0 \\
    0 & -1 & 0 & 0 \\
    0 & 0 & -3 & 0 \\
    0 & 0 & 0 & -4
  \end{pmatrix}$.
  
  Por tanto:
  \[ \PP_t = e^{t\A} = Q e^{tD} Q^{-1} = Q \begin{pmatrix}
    1 & 0 & 0 & 0 \\
    0 & e^{-t} & 0 & 0 \\
    0 & 0 & e^{-3t} & 0 \\
    0 & 0 & 0 & e^{-4t}
  \end{pmatrix} Q^{-1}. \]
  
  La distribuci\'on estacionaria $\pi$ satisface $\pi\A = 0$. Resolviendo el sistema:
  \begin{align*}
    -3\pi_0 + 2\pi_1 + 2\pi_2 + \pi_3 &= 0 \\
    \pi_0 - 4\pi_1 + \pi_2 + \pi_3 &= 0 \\
    \pi_0 + \pi_1 - 4\pi_2 + \pi_3 &= 0 \\
    \pi_0 + \pi_1 + \pi_2 - 3\pi_3 &= 0 \\
    \pi_0 + \pi_1 + \pi_2 + \pi_3 &= 1
  \end{align*}
  
  Por simetr\'ia, $\pi_1 = \pi_2$. Sea $\pi_1 = \pi_2 = a$, $\pi_0 = b$, $\pi_3 = c$. De la segunda ecuaci\'on: $b - 4a + a + c = 0 \Rightarrow b - 3a + c = 0$. De la cuarta: $b + a + a - 3c = 0 \Rightarrow b + 2a - 3c = 0$. Restando: $(b+2a-3c) - (b-3a+c) = 0 \Rightarrow 5a - 4c = 0 \Rightarrow c = 5a/4$. Sustituyendo en $b - 3a + 5a/4 = 0 \Rightarrow b - 12a/4 + 5a/4 = 0 \Rightarrow b = 7a/4$. Normalizaci\'on: $b + a + a + c = 7a/4 + a + a + 5a/4 = (7a+4a+4a+5a)/4 = 20a/4 = 5a = 1 \Rightarrow a = 1/5$. Por tanto:
  \[ \pi = \left(\frac{7}{20}, \frac{1}{5}, \frac{1}{5}, \frac{1}{4}\right) = \left(\frac{7}{20}, \frac{4}{20}, \frac{4}{20}, \frac{5}{20}\right). \]
\end{Ejemplo}

\subsection{Interpretaci\'on mediante Tiempos de Espera Exponenciales}

\begin{Teorema}
  Sea $X_0 = x$ y definamos $T = \inf\{t: X_t \neq x\}$, el tiempo hasta el primer cambio de estado. Entonces $T$ tiene distribuci\'on exponencial con par\'ametro $\alpha(x)$.
\end{Teorema}

\begin{proof}
  Para cualquier $s, t \geq 0$, por la propiedad de Markov y la homogeneidad:
  \[ \Prob\{T > s+t \mid T > s\} = \Prob\{T > t\}. \]
  Esto es la propiedad de p\'erdida de memoria, que caracteriza a la distribuci\'on exponencial. Adem\'as, de \eqref{eq:stay} tenemos:
  \[ \Prob\{T \leq \Delta t\} = 1 - \Prob\{T > \Delta t\} = 1 - (1 - \alpha(x)\Delta t + o(\Delta t)) = \alpha(x)\Delta t + o(\Delta t). \]
  Esto implica que la tasa es $\alpha(x)$. M\'as formalmente, sea $F(t) = \Prob\{T > t\}$. Entonces $F(s+t) = F(s)F(t)$ para todo $s,t \geq 0$, y $F$ es continua por la derecha (de hecho, continua). La \'unica soluci\'on con $F(0)=1$ es $F(t) = e^{-\lambda t}$ para alg\'un $\lambda \geq 0$. De $\Prob\{T \leq \Delta t\} = \lambda \Delta t + o(\Delta t)$, obtenemos $\lambda = \alpha(x)$.
\end{proof}

\begin{Teorema}
  Condicionado a que el primer salto ocurre en tiempo $T$, la probabilidad de que la cadena salte al estado $y \neq x$ es
  \[ \Prob\{X_T = y\} = \frac{\alpha(x,y)}{\alpha(x)}. \]
\end{Teorema}

\begin{proof}
  Consideramos un intervalo peque\~no $[t, t+\Delta t]$ y aplicamos \eqref{eq:jump}:
  \[ \Prob\{X_{t+\Delta t} = y, X_s = x \text{ para } s \in [0,t] \mid X_0 = x\} = e^{-\alpha(x)t} \cdot \alpha(x,y)\Delta t + o(\Delta t). \]
  Esto se debe a que la probabilidad de no haber saltado hasta $t$ es $e^{-\alpha(x)t}$, y en el intervalo $[t,t+\Delta t]$ la probabilidad de saltar a $y$ es aproximadamente $\alpha(x,y)\Delta t$, con independencia aproximada.

  La probabilidad de que el primer salto ocurra en $[t, t+\Delta t]$ (sin especificar destino) es:
  \[ \Prob\{T \in [t,t+\Delta t]\} = e^{-\alpha(x)t}\alpha(x)\Delta t + o(\Delta t). \]

  Por tanto, la probabilidad condicionada es:
  \[ \Prob\{X_T = y \mid T \in [t,t+\Delta t]\} = \frac{e^{-\alpha(x)t}\alpha(x,y)\Delta t + o(\Delta t)}{e^{-\alpha(x)t}\alpha(x)\Delta t + o(\Delta t)} \xrightarrow{\Delta t \to 0} \frac{\alpha(x,y)}{\alpha(x)}. \]
  Esto es independiente de $t$, como debe ser.
\end{proof}

\begin{Observacion}[Interpretaci\'on de los relojes]
  Podemos pensar que en cada estado $y \neq x$ colocamos un "reloj" independiente que suena en un tiempo exponencial con tasa $\alpha(x,y)$. La cadena permanece en $x$ hasta que el primer reloj suena, y entonces se mueve al estado correspondiente. Esta interpretaci\'on es consistente con los teoremas anteriores.
\end{Observacion}

\subsection{Comportamiento Asint\'otico y Distribuci\'on Estacionaria}

\begin{Definicion}
  Una cadena en tiempo continuo es \textbf{irreducible} si todos los estados se comunican, es decir, para cada $x, y \in S$ existe una secuencia $z_1, \ldots, z_j$ tal que
  \[ \alpha(x,z_1), \alpha(z_1,z_2), \ldots, \alpha(z_{j-1},z_j), \alpha(z_j,y) > 0. \]
\end{Definicion}

\begin{Lema}
  Sea $\A$ el generador infinitesimal de una cadena irreducible en tiempo continuo con espacio de estados finito. Sea $a > \max_x \alpha(x)$. Entonces la matriz $P = \frac{1}{a}\A + \I$ es una matriz estoc\'astica (sus filas suman 1 y todas sus entradas son no negativas) irreducible y aperi\'odica.
\end{Lema}

\begin{proof}
  Para cualquier $x$, la suma de la fila $x$ de $P$ es:
  \[ \sum_y P(x,y) = \frac{1}{a}\sum_y \A(x,y) + \sum_y \delta_{xy} = \frac{1}{a}\cdot 0 + 1 = 1. \]
  Las entradas no diagonales son $P(x,y) = \frac{1}{a}\alpha(x,y) \geq 0$ para $x \neq y$. La entrada diagonal es $P(x,x) = 1 - \frac{\alpha(x)}{a} \geq 0$ porque $a > \alpha(x)$. Por tanto, $P$ es estoc\'astica.

  La irreducibilidad se conserva porque $\alpha(x,y) > 0$ implica $P(x,y) > 0$. Adem\'as, $P(x,x) > 0$ para alg\'un $x$ (de hecho, para todos porque $a > \alpha(x)$), lo que implica que la cadena en tiempo discreto es aperi\'odica.
\end{proof}

\begin{Teorema}
  Para una cadena irreducible en tiempo continuo con espacio de estados finito, existe una \'unica distribuci\'on de probabilidad $\pi$ tal que
  \[ \lim_{t\to\infty} \PP_t = \ind\pi, \]
  donde $\ind\pi$ es la matriz con todas las filas iguales a $\pi$. La distribuci\'on $\pi$ satisface
  \[ \pi \A = \mathbf{0}. \]
  Adem\'as, todos los dem\'as eigenvalores de $\A$ tienen parte real negativa.
\end{Teorema}

\begin{proof}
  Elegimos $a > \max_x \alpha(x)$ y definimos $P = \frac{1}{a}\A + \I$. Por el Lema, $P$ es una matriz estoc\'astica irreducible y aperi\'odica. Por la teor\'ia de cadenas de Markov en tiempo discreto, existe una \'unica distribuci\'on estacionaria $\pi$ tal que $\pi P = \pi$ y $\lim_{n\to\infty} P^n = \ind\pi$.

  De $P = \frac{1}{a}\A + \I$, obtenemos:
  \[ \pi P = \pi \Rightarrow \pi\left(\frac{1}{a}\A + \I\right) = \pi \Rightarrow \frac{1}{a}\pi\A + \pi = \pi \Rightarrow \pi\A = 0. \]

  Ahora, para la convergencia en tiempo continuo, observamos que $\PP_t = e^{t\A}$. Podemos escribir $\A$ en t\'erminos de $P$: $\A = a(P - \I)$. Entonces:
  \[ \PP_t = e^{t a(P - \I)} = e^{-a t} e^{a t P}. \]

  Descomponemos $P$ en sus eigenvalores. Como $P$ es estoc\'astica irreducible y aperi\'odica, tiene eigenvalor 1 con eigenvector derecho $\ind$ y eigenvector izquierdo $\pi$. Los dem\'as eigenvalores $\mu$ satisfacen $|\mu| < 1$. Por tanto, los eigenvalores de $\A$ son $\lambda = a(\mu - 1)$. Para $\mu = 1$, obtenemos $\lambda = 0$. Para $|\mu| < 1$, tenemos $\operatorname{Re}(\lambda) = a(\operatorname{Re}(\mu) - 1) < 0$ porque $\operatorname{Re}(\mu) < 1$.

  La descomposici\'on espectral nos da:
  \[ \PP_t = \ind\pi + \sum_{\lambda \neq 0} e^{\lambda t} \text{(t\'erminos de proyecci\'on)}. \]
  Como $\operatorname{Re}(\lambda) < 0$ para $\lambda \neq 0$, estos t\'erminos tienden a 0 cuando $t \to \infty$. Por tanto, $\lim_{t\to\infty} \PP_t = \ind\pi$.
\end{proof}

\begin{Observacion}
  En tiempo continuo no existe el fen\'omeno de periodicidad. Para cualquier cadena irreducible en tiempo continuo, $\PP_t$ tiene todas sus entradas estrictamente positivas para todo $t > 0$.
\end{Observacion}

\begin{proof}
  Por irreducibilidad, para cualquier par de estados $i, j$, existe un camino $i = x_0, x_1, \ldots, x_n = j$ con tasas positivas $\alpha(x_k, x_{k+1}) > 0$. Podemos elegir tiempos $0 < t_1 < t_2 < \cdots < t_n = t$. La probabilidad de seguir este camino con saltos exactamente en esos tiempos es:
  \[ e^{-\alpha(x_0)t_1} \cdot \alpha(x_0,x_1) \cdot e^{-\alpha(x_1)(t_2-t_1)} \cdot \alpha(x_1,x_2) \cdots e^{-\alpha(x_{n-1})(t_n-t_{n-1})} > 0. \]
  Integrando sobre todas las posibles secuencias de tiempos de salto que suman $t$, obtenemos $p_t(i,j) > 0$.
\end{proof}

\subsection{Tiempos Medios de Paso}

\begin{Definicion}
  Sea $z$ un estado fijo. Definimos el \textbf{tiempo medio de paso} a $z$ comenzando en $x$ como
  \[ b(x) = \E(Y \mid X_0 = x), \quad \text{donde } Y = \inf\{t: X_t = z\}. \]
  Claramente $b(z) = 0$.
\end{Definicion}

\begin{Teorema}
  Para $x \neq z$, la funci\'on $b(x)$ satisface
  \[ \alpha(x)b(x) = 1 + \sum_{y \neq x, z} \alpha(x,y) b(y). \]
\end{Teorema}

\begin{proof}
  Condicionamos en el primer tiempo de salto $T$ desde $x$. Sea $Y$ el tiempo de primera llegada a $z$. Entonces:
  \[ Y = T + Y' \]
  donde $Y'$ es el tiempo restante despu\'es del primer salto. Por la propiedad de Markov, condicionado en $X_T = y$, $Y'$ tiene la misma distribuci\'on que $Y$ comenzando en $y$ (por la homogeneidad temporal). Por tanto:
  \[ \E(Y \mid X_0 = x) = \E(T \mid X_0 = x) + \sum_{y \in S} \Prob\{X_T = y \mid X_0 = x\} \E(Y \mid X_0 = y). \]

  Sabemos que $\E(T \mid X_0 = x) = 1/\alpha(x)$ porque $T \sim \operatorname{Exp}(\alpha(x))$. Tambi\'en $\Prob\{X_T = y \mid X_0 = x\} = \alpha(x,y)/\alpha(x)$ para $y \neq x$, y $\Prob\{X_T = x \mid X_0 = x\} = 0$ (casi seguramente la cadena cambia de estado).

  Sustituyendo:
  \[ b(x) = \frac{1}{\alpha(x)} + \sum_{y \neq x} \frac{\alpha(x,y)}{\alpha(x)} b(y). \]

  Separamos la suma en $y = z$ y $y \neq z$:
  \[ b(x) = \frac{1}{\alpha(x)} + \frac{\alpha(x,z)}{\alpha(x)} b(z) + \sum_{y \neq x, z} \frac{\alpha(x,y)}{\alpha(x)} b(y). \]

  Como $b(z) = 0$, obtenemos:
  \[ b(x) = \frac{1}{\alpha(x)} + \sum_{y \neq x, z} \frac{\alpha(x,y)}{\alpha(x)} b(y). \]

  Multiplicando ambos lados por $\alpha(x)$:
  \[ \alpha(x) b(x) = 1 + \sum_{y \neq x, z} \alpha(x,y) b(y). \]
\end{proof}

\begin{Corolario}
  Sea $\tilde{\A}$ la matriz obtenida de $\A$ eliminando la fila y columna correspondientes al estado $z$. Sea $\bar{b}$ el vector de tiempos medios de paso para los estados $x \neq z$, y $\bar{1}$ el vector de unos. Entonces
  \[ \tilde{\A} \bar{b} = -\bar{1}, \]
  y por tanto
  \[ \bar{b} = [-\tilde{\A}]^{-1} \bar{1}. \]
\end{Corolario}

\begin{proof}
  Para $x \neq z$, la ecuaci\'on $\alpha(x) b(x) = 1 + \sum_{y \neq x, z} \alpha(x,y) b(y)$ se puede reescribir como:
  \[ -\alpha(x) b(x) + \sum_{y \neq x, z} \alpha(x,y) b(y) = -1. \]

  El lado izquierdo es precisamente la componente $x$ de $\tilde{\A} \bar{b}$, porque $\tilde{\A}(x,x) = -\alpha(x)$ y $\tilde{\A}(x,y) = \alpha(x,y)$ para $y \neq x$. Por tanto:
  \[ (\tilde{\A} \bar{b})_x = -1 \quad \text{para todo } x \neq z. \]

  Esto es $\tilde{\A} \bar{b} = -\bar{1}$. Como la cadena es irreducible, $\tilde{\A}$ es invertible (todos sus eigenvalores tienen parte real negativa), por lo que $\bar{b} = -\tilde{\A}^{-1}\bar{1} = [-\tilde{\A}]^{-1}\bar{1}$.
\end{proof}

\begin{Ejemplo}[C\'alculo de tiempo medio de paso]
  Para la cadena del Ejemplo anterior (la de cuatro estados), calculemos el tiempo esperado para llegar del estado 0 al estado 3. Tenemos $z = 3$, y la matriz $\tilde{\A}$ es:
  \[ \tilde{\A} = \begin{pmatrix}
    -3 & 1 & 1 \\
    2 & -4 & 1 \\
    2 & 1 & -4
  \end{pmatrix}. \]

  Calculamos $-\tilde{\A}$ y resolvemos $-\tilde{\A} \bar{b} = \bar{1}$:
  \[ \begin{cases}
    3b_0 - b_1 - b_2 = 1 \\
    -2b_0 + 4b_1 - b_2 = 1 \\
    -2b_0 - b_1 + 4b_2 = 1
  \end{cases} \]

  De la primera ecuaci\'on: $b_1 = 3b_0 - b_2 - 1$.
  Sustituyendo en la segunda: $-2b_0 + 4(3b_0 - b_2 - 1) - b_2 = 1 \Rightarrow -2b_0 + 12b_0 - 4b_2 - 4 - b_2 = 1 \Rightarrow 10b_0 - 5b_2 = 5 \Rightarrow 2b_0 - b_2 = 1 \Rightarrow b_2 = 2b_0 - 1$.

  Sustituyendo en la tercera: $-2b_0 - (2b_0 - 1) + 4(2b_0 - 1) = 1 \Rightarrow -2b_0 - 2b_0 + 1 + 8b_0 - 4 = 1 \Rightarrow 4b_0 - 3 = 1 \Rightarrow 4b_0 = 4 \Rightarrow b_0 = 1$.

  Entonces $b_2 = 2\cdot 1 - 1 = 1$, y $b_1 = 3\cdot 1 - 1 - 1 = 1$.

  Por tanto, $\bar{b} = (1,1,1)^T$. El tiempo esperado del estado 0 al estado 3 es 1.
\end{Ejemplo}

\section{Procesos de Nacimiento y Muerte}

\begin{Definicion}
  Un \textbf{proceso de nacimiento y muerte} es una cadena de Markov en tiempo continuo con espacio de estados $\{0,1,2,\ldots\}$ donde los \'unicos cambios posibles son de $n$ a $n+1$ (nacimiento) o de $n$ a $n-1$ (muerte). Definimos:
  \begin{itemize}
    \item Tasas de nacimiento: $\lambda_n \geq 0$, $n = 0,1,2,\ldots$
    \item Tasas de muerte: $\mu_n \geq 0$, $n = 1,2,3,\ldots$, con $\mu_0 = 0$.
  \end{itemize}
  Las probabilidades infinitesimales son:
  \begin{align*}
    \Prob\{X_{t+\Delta t} = n \mid X_t = n\} &= 1 - (\lambda_n + \mu_n)\Delta t + o(\Delta t),\\
    \Prob\{X_{t+\Delta t} = n+1 \mid X_t = n\} &= \lambda_n \Delta t + o(\Delta t),\\
    \Prob\{X_{t+\Delta t} = n-1 \mid X_t = n\} &= \mu_n \Delta t + o(\Delta t).
  \end{align*}
\end{Definicion}

\begin{Proposicion}
  Las probabilidades $P_n(t) = \Prob\{X_t = n\}$ satisfacen las ecuaciones diferenciales:
  \[ P_n'(t) = \mu_{n+1} P_{n+1}(t) + \lambda_{n-1} P_{n-1}(t) - (\mu_n + \lambda_n) P_n(t), \quad n \geq 0, \]
  con el convenio $\lambda_{-1} = 0$ y $\mu_0 = 0$.
\end{Proposicion}

\begin{proof}
  Aplicamos la misma t\'ecnica que para cadenas generales, pero ahora las transiciones son solo a estados vecinos:
  \[
  P_n(t+\Delta t) = P_n(t)(1 - (\lambda_n + \mu_n)\Delta t + o(\Delta t)) + P_{n-1}(t)(\lambda_{n-1}\Delta t + o(\Delta t)) + P_{n+1}(t)(\mu_{n+1}\Delta t + o(\Delta t)).
  \]
  Reordenando:
  \[
  \frac{P_n(t+\Delta t) - P_n(t)}{\Delta t} = -(\lambda_n + \mu_n)P_n(t) + \lambda_{n-1}P_{n-1}(t) + \mu_{n+1}P_{n+1}(t) + \frac{o(\Delta t)}{\Delta t}(P_n(t) + P_{n-1}(t) + P_{n+1}(t)).
  \]
  Tomando l\'imite $\Delta t \to 0$, obtenemos la ecuaci\'on.
\end{proof}

\subsection{Ejemplos de Procesos de Nacimiento y Muerte}

\begin{Ejemplo}[Proceso de Poisson]
  El proceso de Poisson con tasa $\lambda$ es un proceso de nacimiento y muerte con $\lambda_n = \lambda$, $\mu_n = 0$ para todo $n$.
\end{Ejemplo}

\begin{Ejemplo}[Modelos de colas M/M]
  Sea $X_t$ el n\'umero de personas en una cola.
  \begin{description}
    \item[Cola M/M/1:] Un servidor. $\lambda_n = \lambda$, $\mu_n = \mu$ para $n \geq 1$, $\mu_0 = 0$.
    \item[Cola M/M/k:] $k$ servidores. $\lambda_n = \lambda$, $\mu_n = \min(n,k)\mu$.
    \item[Cola M/M/$\infty$:] Infinitos servidores. $\lambda_n = \lambda$, $\mu_n = n\mu$.
  \end{description}
  La notaci\'on M/M indica que tanto los tiempos entre llegadas como los tiempos de servicio son exponenciales (Markovianos).
\end{Ejemplo}

\begin{Ejemplo}[Modelo de poblaci\'on]
  Cada individuo produce nuevos individuos a tasa $\lambda$ y muere a tasa $\mu$, independientemente. Esto da $\lambda_n = n\lambda$, $\mu_n = n\mu$. Cuando $\mu = 0$, se llama proceso de Yule.
\end{Ejemplo}

\begin{Ejemplo}[Modelo de poblaci\'on con inmigraci\'on]
  Adem\'as de los nacimientos y muertes internos, llegan nuevos individuos del exterior a tasa $\nu$. Entonces $\lambda_n = n\lambda + \nu$, $\mu_n = n\mu$.
\end{Ejemplo}

\begin{Ejemplo}[Modelo de crecimiento r\'apido]
  $\lambda_n = n^2\lambda$, $\mu_n = 0$. Este modelo puede explotar en tiempo finito.
\end{Ejemplo}

\subsection{Clasificaci\'on de Procesos de Nacimiento y Muerte}

\begin{Definicion}
  Un proceso de nacimiento y muerte es \textbf{irreducible} si todos los estados se comunican, lo cual ocurre si y solo si $\lambda_n > 0$ para todo $n$ y $\mu_n > 0$ para todo $n \geq 1$.
\end{Definicion}

\begin{Definicion}
  Una cadena irreducible es \textbf{recurrente} si con probabilidad 1 se regresa a cada estado; en caso contrario es \textbf{transitoria}. Si es recurrente y existe una distribuci\'on estacionaria, es \textbf{positivamente recurrente}; si no, es \textbf{nulamente recurrente}.
\end{Definicion}

\begin{Observacion}
  Para procesos de nacimiento y muerte, existe una cadena en tiempo discreto embebida que sigue la cadena "cuando se mueve". Sus probabilidades de transici\'on son:
  \[ p(n,n-1) = \frac{\mu_n}{\mu_n + \lambda_n}, \quad p(n,n+1) = \frac{\lambda_n}{\mu_n + \lambda_n}. \]
  La cadena en tiempo continuo es recurrente si y solo si la cadena embebida es recurrente.
\end{Observacion}

\begin{Teorema}[Criterio de transiencia]
  El proceso de nacimiento y muerte es transitorio si y solo si
  \[ \sum_{n=1}^{\infty} \frac{\mu_1 \mu_2 \cdots \mu_n}{\lambda_1 \lambda_2 \cdots \lambda_n} < \infty. \]
\end{Teorema}

\begin{proof}
  Sea $a(n)$ la probabilidad de que, comenzando en $n$, la cadena alcance alguna vez el estado 0. Entonces $a(0)=1$. Para $n \geq 1$, condicionando en el primer paso:
  \[ a(n) = \frac{\mu_n}{\mu_n + \lambda_n} a(n-1) + \frac{\lambda_n}{\mu_n + \lambda_n} a(n+1). \]
  Multiplicando por $\mu_n + \lambda_n$:
  \[ (\mu_n + \lambda_n) a(n) = \mu_n a(n-1) + \lambda_n a(n+1). \]
  Reordenando:
  \[ \lambda_n (a(n+1) - a(n)) = \mu_n (a(n) - a(n-1)). \]
  Definimos $d_n = a(n) - a(n+1)$. Entonces:
  \[ \lambda_n d_n = \mu_n d_{n-1} \Rightarrow d_n = \frac{\mu_n}{\lambda_n} d_{n-1}. \]
  Iterando:
  \[ d_n = \frac{\mu_1 \mu_2 \cdots \mu_n}{\lambda_1 \lambda_2 \cdots \lambda_n} d_0. \]
  Como $a(0)=1$, tenemos $d_0 = 1 - a(1)$. Entonces:
  \[ a(n+1) = a(0) - \sum_{j=0}^n d_j = 1 - (1 - a(1)) \sum_{j=0}^n \frac{\mu_1 \cdots \mu_j}{\lambda_1 \cdots \lambda_j}, \]
  con el convenio de que el t\'ermino $j=0$ es 1.

  Si la cadena es transitoria, entonces $a(n) \to 0$ cuando $n \to \infty$. Esto requiere que la suma $\sum_{j=0}^\infty \frac{\mu_1 \cdots \mu_j}{\lambda_1 \cdots \lambda_j}$ converja, y entonces $a(1) = 1 - 1/\sum_{j=0}^\infty \frac{\mu_1 \cdots \mu_j}{\lambda_1 \cdots \lambda_j}$. Por tanto, la condici\'on de convergencia es equivalente a la transiencia.

  Si la suma diverge, entonces la \'unica soluci\'on con $a(0)=1$ y $0 \leq a(n) \leq 1$ es $a(n)=1$ para todo $n$, lo que corresponde a recurrencia.
\end{proof}

\begin{Teorema}[Criterio de recurrencia positiva]
  El proceso de nacimiento y muerte es positivamente recurrente si y solo si
  \[ \sum_{n=0}^{\infty} \frac{\lambda_0 \lambda_1 \cdots \lambda_{n-1}}{\mu_1 \mu_2 \cdots \mu_n} < \infty, \]
  donde el t\'ermino $n=0$ se interpreta como 1. En este caso, la distribuci\'on estacionaria es
  \[ \pi(n) = \frac{\lambda_0 \lambda_1 \cdots \lambda_{n-1}}{\mu_1 \mu_2 \cdots \mu_n} \cdot \frac{1}{C}, \quad \text{con } C = \sum_{n=0}^\infty \frac{\lambda_0 \lambda_1 \cdots \lambda_{n-1}}{\mu_1 \mu_2 \cdots \mu_n}. \]
\end{Teorema}

\begin{proof}
  Buscamos una distribuci\'on estacionaria $\pi$ que satisfaga las ecuaciones de balance:
  \[ 0 = \lambda_{n-1} \pi(n-1) + \mu_{n+1} \pi(n+1) - (\lambda_n + \mu_n) \pi(n), \quad n \geq 0. \]
  Para $n=0$, como $\mu_0 = 0$ y $\lambda_{-1}=0$, obtenemos:
  \[ 0 = \mu_1 \pi(1) - \lambda_0 \pi(0) \Rightarrow \pi(1) = \frac{\lambda_0}{\mu_1} \pi(0). \]
  Para $n \geq 1$, podemos reescribir la ecuaci\'on como:
  \[ \mu_{n+1} \pi(n+1) - \lambda_n \pi(n) = \mu_n \pi(n) - \lambda_{n-1} \pi(n-1). \]
  Esto muestra que la cantidad $\mu_n \pi(n) - \lambda_{n-1} \pi(n-1)$ es constante en $n$. Evaluando en $n=1$:
  \[ \mu_1 \pi(1) - \lambda_0 \pi(0) = \mu_1 \cdot \frac{\lambda_0}{\mu_1} \pi(0) - \lambda_0 \pi(0) = 0. \]
  Por tanto, la constante es 0, y tenemos:
  \[ \mu_{n+1} \pi(n+1) = \lambda_n \pi(n) \Rightarrow \pi(n+1) = \frac{\lambda_n}{\mu_{n+1}} \pi(n). \]
  Iterando:
  \[ \pi(n) = \frac{\lambda_0 \lambda_1 \cdots \lambda_{n-1}}{\mu_1 \mu_2 \cdots \mu_n} \pi(0). \]
  La condici\'on de normalizaci\'on $\sum_{n=0}^\infty \pi(n) = 1$ requiere que la serie $\sum_{n=0}^\infty \frac{\lambda_0 \lambda_1 \cdots \lambda_{n-1}}{\mu_1 \mu_2 \cdots \mu_n}$ converja. Si converge, definimos $C$ como su suma, y entonces $\pi(0) = 1/C$, obteniendo la f\'ormula.
\end{proof}

\subsection{Aplicaciones a Modelos de Colas}

\begin{Ejemplo}[Cola M/M/1]
  Para la cola M/M/1, $\lambda_n = \lambda$, $\mu_n = \mu$ para $n \geq 1$, $\mu_0 = 0$.
  
  Criterio de transiencia:
  \[ \sum_{n=1}^\infty \frac{\mu_1 \cdots \mu_n}{\lambda_1 \cdots \lambda_n} = \sum_{n=1}^\infty \left(\frac{\mu}{\lambda}\right)^n. \]
  Esta serie converge si $\mu/\lambda < 1$, es decir, si $\mu < \lambda$. Por tanto, la cola es transitoria si $\mu < \lambda$ (la tasa de servicio es menor que la tasa de llegada), lo que significa que la cola tiende a crecer indefinidamente.
  
  Criterio de recurrencia positiva:
  \[ \sum_{n=0}^\infty \frac{\lambda_0 \cdots \lambda_{n-1}}{\mu_1 \cdots \mu_n} = \sum_{n=0}^\infty \left(\frac{\lambda}{\mu}\right)^n = \frac{1}{1 - \lambda/\mu} \quad \text{si } \lambda < \mu. \]
  La serie converge si y solo si $\lambda < \mu$. En este caso, la distribuci\'on estacionaria es:
  \[ \pi(n) = \left(1 - \frac{\lambda}{\mu}\right) \left(\frac{\lambda}{\mu}\right)^n, \quad n \geq 0. \]
  Esta es una distribuci\'on geom\'etrica. La longitud esperada de la cola en equilibrio es:
  \[ \E[X] = \sum_{n=0}^\infty n \pi(n) = \frac{\lambda/\mu}{1 - \lambda/\mu} = \frac{\lambda}{\mu - \lambda}. \]
  Observamos que esta esperanza tiende a infinito cuando $\lambda \to \mu^-$.
\end{Ejemplo}

\begin{Ejemplo}[Cola M/M/k]
  Para la cola M/M/k, $\lambda_n = \lambda$, $\mu_n = \min(n,k)\mu$. Para $n \geq k$, $\mu_n = k\mu$.
  
  Para $n \geq k$, tenemos:
  \[ \frac{\mu_1 \cdots \mu_n}{\lambda_1 \cdots \lambda_n} = \frac{[\mu \cdot 2\mu \cdots (k-1)\mu] \cdot (k\mu)^{n-k+1}}{\lambda^n} = \frac{(k-1)! k^{n-k+1} \mu^n}{\lambda^n}. \]
  Esto es $\frac{(k-1)!}{k^{k-1}} \left(\frac{k\mu}{\lambda}\right)^n \cdot \text{constante}$. La serie $\sum_{n=k}^\infty \left(\frac{k\mu}{\lambda}\right)^n$ converge si y solo si $k\mu < \lambda$. Por tanto, la cola es transitoria si $k\mu < \lambda$.
  
  Para la recurrencia positiva, la serie estacionaria converge si y solo si $\lambda < k\mu$. En este caso, la distribuci\'on estacionaria existe pero su forma es m\'as complicada.
\end{Ejemplo}

\begin{Ejemplo}[Cola M/M/$\infty$]
  Para la cola M/M/$\infty$, $\lambda_n = \lambda$, $\mu_n = n\mu$.
  
  La serie para la transiencia es $\sum_{n=1}^\infty \frac{n! \mu^n}{\lambda^n}$, que diverge para todo $\lambda, \mu > 0$ (por la f\'ormula de Stirling, $n! \sim \sqrt{2\pi n}(n/e)^n$, y el t\'ermino $n^n$ domina). Por tanto, la cola es siempre recurrente.
  
  La serie para la recurrencia positiva es:
  \[ \sum_{n=0}^\infty \frac{\lambda^n}{n! \mu^n} = \sum_{n=0}^\infty \frac{(\lambda/\mu)^n}{n!} = e^{\lambda/\mu} < \infty. \]
  Por tanto, la cola es siempre positivamente recurrente. La distribuci\'on estacionaria es:
  \[ \pi(n) = e^{-\lambda/\mu} \frac{(\lambda/\mu)^n}{n!}, \]
  una distribuci\'on de Poisson con par\'ametro $\lambda/\mu$. La longitud esperada de la cola en equilibrio es $\lambda/\mu$.
\end{Ejemplo}

\subsection{Procesos de Nacimiento Puro}

\begin{Definicion}
  Un proceso de nacimiento y muerte con $\mu_n = 0$ para todo $n$ se llama \textbf{proceso de nacimiento puro}.
\end{Definicion}

\begin{Ejemplo}[Proceso de Yule]
  El proceso de Yule tiene $\lambda_n = n\lambda$, $\mu_n = 0$. Supongamos que $X_0 = 1$.
  
  Las ecuaciones diferenciales para $P_n(t) = \Prob\{X_t = n\}$ son:
  \[ P_n'(t) = (n-1)\lambda P_{n-1}(t) - n\lambda P_n(t), \quad n \geq 1, \]
  con $P_1(0) = 1$.
  
  Resolvemos por inducci\'on. Para $n=1$:
  \[ P_1'(t) = -\lambda P_1(t) \Rightarrow P_1(t) = e^{-\lambda t}. \]
  Para $n=2$:
  \[ P_2'(t) = \lambda P_1(t) - 2\lambda P_2(t) = \lambda e^{-\lambda t} - 2\lambda P_2(t). \]
  Esta es una ecuaci\'on lineal de primer orden de la forma $P_2' + 2\lambda P_2 = \lambda e^{-\lambda t}$. Multiplicamos por $e^{2\lambda t}$:
  \[ \frac{d}{dt}(P_2 e^{2\lambda t}) = \lambda e^{-\lambda t} e^{2\lambda t} = \lambda e^{\lambda t}. \]
  Integrando: $P_2 e^{2\lambda t} = \int \lambda e^{\lambda t} dt = e^{\lambda t} + C$. Con $P_2(0)=0$, obtenemos $0 = 1 + C \Rightarrow C = -1$. Por tanto:
  \[ P_2 e^{2\lambda t} = e^{\lambda t} - 1 \Rightarrow P_2(t) = e^{-\lambda t} - e^{-2\lambda t} = e^{-\lambda t}(1 - e^{-\lambda t}). \]
  
  Para $n=3$:
  \[ P_3'(t) = 2\lambda P_2(t) - 3\lambda P_3(t) = 2\lambda e^{-\lambda t}(1 - e^{-\lambda t}) - 3\lambda P_3(t). \]
  Multiplicamos por $e^{3\lambda t}$:
  \[ \frac{d}{dt}(P_3 e^{3\lambda t}) = 2\lambda e^{-\lambda t}(1 - e^{-\lambda t}) e^{3\lambda t} = 2\lambda e^{2\lambda t}(1 - e^{-\lambda t}) = 2\lambda(e^{2\lambda t} - e^{\lambda t}). \]
  Integrando: $P_3 e^{3\lambda t} = 2\lambda \int (e^{2\lambda t} - e^{\lambda t}) dt = 2\lambda \left(\frac{e^{2\lambda t}}{2\lambda} - \frac{e^{\lambda t}}{\lambda}\right) + C = e^{2\lambda t} - 2e^{\lambda t} + C$.
  Con $P_3(0)=0$, tenemos $0 = 1 - 2 + C \Rightarrow C = 1$. Por tanto:
  \[ P_3 e^{3\lambda t} = e^{2\lambda t} - 2e^{\lambda t} + 1 = (e^{\lambda t} - 1)^2 \Rightarrow P_3(t) = e^{-3\lambda t}(e^{\lambda t} - 1)^2 = e^{-\lambda t}(1 - e^{-\lambda t})^2. \]
  
  Por inducci\'on, se puede demostrar que:
  \[ P_n(t) = e^{-\lambda t}(1 - e^{-\lambda t})^{n-1}, \quad n \geq 1. \]
  Esta es una distribuci\'on geom\'etrica con par\'ametro $e^{-\lambda t}$.
  
  La esperanza de $X_t$ es:
  \[ \E[X_t] = \sum_{n=1}^\infty n e^{-\lambda t}(1 - e^{-\lambda t})^{n-1} = e^{\lambda t}. \]
  Podemos obtener esto tambi\'en de una ecuaci\'on diferencial para $m(t) = \E[X_t]$:
  \[ m'(t) = \sum_{n=1}^\infty n P_n'(t) = \sum_{n=1}^\infty n[(n-1)\lambda P_{n-1} - n\lambda P_n] = \lambda \sum_{n=1}^\infty (n-1)P_{n-1} - \lambda \sum_{n=1}^\infty n^2 P_n. \]
  El primer t\'ermino: $\lambda \sum_{n=1}^\infty (n-1)P_{n-1} = \lambda \sum_{k=0}^\infty k P_k = \lambda m(t)$. El segundo t\'ermino: $-\lambda \sum_{n=1}^\infty n^2 P_n = -\lambda \E[X_t^2]$. Esto no parece simplificar directamente. Mejor usamos otro m\'etodo:
  \[ m'(t) = \sum_{n=1}^\infty n P_n'(t) = \sum_{n=1}^\infty n[(n-1)\lambda P_{n-1} - n\lambda P_n] = \lambda \sum_{n=1}^\infty n(n-1)P_{n-1} - \lambda \sum_{n=1}^\infty n^2 P_n. \]
  En el primer t\'ermino, hacemos $k = n-1$: $\sum_{n=1}^\infty n(n-1)P_{n-1} = \sum_{k=0}^\infty (k+1)k P_k = \E[X_t^2] + \E[X_t]$. El segundo t\'ermino es $\E[X_t^2]$. Por tanto:
  \[ m'(t) = \lambda (\E[X_t^2] + \E[X_t] - \E[X_t^2]) = \lambda \E[X_t] = \lambda m(t). \]
  Con $m(0)=1$, obtenemos $m(t) = e^{\lambda t}$, correcto.
  
  Otra perspectiva: sea $Y_n$ el tiempo en que la poblaci\'on alcanza $n$ por primera vez. Entonces $Y_n = T_1 + \cdots + T_{n-1}$, donde $T_i$ es el tiempo entre la llegada del individuo $i$ y el $i+1$. Los $T_i$ son independientes con $T_i \sim \operatorname{Exp}(i\lambda)$. Por tanto:
  \[ \E[Y_n] = \sum_{i=1}^{n-1} \frac{1}{i\lambda} \sim \frac{\ln n}{\lambda}, \]
  y $\operatorname{Var}(Y_n) = \sum_{i=1}^{n-1} \frac{1}{i^2\lambda^2} < \infty$. Esto significa que $Y_n$ es aproximadamente $\ln n/\lambda$ con un error acotado.
\end{Ejemplo}

\begin{Ejemplo}[Modelo de crecimiento r\'apido]
  Consideremos $\lambda_n = n^2\lambda$, $\mu_n = 0$. Sea $Y_n$ el tiempo hasta alcanzar $n$ individuos, comenzando con 1. Entonces $Y_n = T_1 + \cdots + T_{n-1}$ con $T_i \sim \operatorname{Exp}(i^2\lambda)$. Entonces:
  \[ \E[Y_\infty] = \sum_{i=1}^\infty \frac{1}{i^2\lambda} = \frac{1}{\lambda} \cdot \frac{\pi^2}{6} < \infty. \]
  Esto implica que con probabilidad 1, $Y_\infty < \infty$, es decir, la poblaci\'on se vuelve infinita en tiempo finito. Este fen\'omeno se llama \textbf{explosi\'on}.
\end{Ejemplo}

\section{Ecuaciones Forward y Backward}

\begin{Definicion}[Ecuaciones forward]
  Las ecuaciones forward se obtienen condicionando en el estado en tiempo $t$:
  \[ p_{t+\Delta t}(x,y) = \sum_z p_t(x,z) p_{\Delta t}(z,y). \]
  Usando las probabilidades infinitesimales:
  \begin{align*}
    p_{t+\Delta t}(x,y) &= p_t(x,y)(1 - \alpha(y)\Delta t + o_y(\Delta t)) + \sum_{z \neq y} p_t(x,z)(\alpha(z,y)\Delta t + o_z(\Delta t)).
  \end{align*}
  Si podemos intercambiar l\'imites y justificar que $\sum_z p_t(x,z) o_z(\Delta t) = o(\Delta t)$, entonces:
  \[ \frac{p_{t+\Delta t}(x,y) - p_t(x,y)}{\Delta t} = -\alpha(y)p_t(x,y) + \sum_{z \neq y} \alpha(z,y)p_t(x,z) + \frac{o(\Delta t)}{\Delta t}. \]
  Tomando $\Delta t \to 0$:
  \[ \frac{\partial}{\partial t} p_t(x,y) = -\alpha(y)p_t(x,y) + \sum_{z \neq y} \alpha(z,y)p_t(x,z). \]
  En forma matricial: $\frac{d}{dt} \PP_t = \PP_t \A$.
\end{Definicion}

\begin{Definicion}[Ecuaciones backward]
  Las ecuaciones backward se obtienen condicionando en el primer intervalo:
  \[ p_{t+\Delta t}(x,y) = \sum_z p_{\Delta t}(x,z) p_t(z,y). \]
  Usando las probabilidades infinitesimales:
  \begin{align*}
    p_{t+\Delta t}(x,y) &= (1 - \alpha(x)\Delta t + o_x(\Delta t))p_t(x,y) + \sum_{z \neq x} (\alpha(x,z)\Delta t + o_x(\Delta t)) p_t(z,y).
  \end{align*}
  Reordenando:
  \[ \frac{p_{t+\Delta t}(x,y) - p_t(x,y)}{\Delta t} = -\alpha(x)p_t(x,y) + \sum_{z \neq x} \alpha(x,z)p_t(z,y) + \frac{o_x(\Delta t)}{\Delta t}\left(p_t(x,y) + \sum_{z \neq x} p_t(z,y)\right). \]
  Tomando $\Delta t \to 0$:
  \[ \frac{\partial}{\partial t} p_t(x,y) = -\alpha(x)p_t(x,y) + \sum_{z \neq x} \alpha(x,z)p_t(z,y). \]
  En forma matricial: $\frac{d}{dt} \PP_t = \A \PP_t$.
\end{Definicion}

\begin{Observacion}
  Para espacios de estados finitos, ambas ecuaciones tienen la misma soluci\'on $\PP_t = e^{t\A}$. Las ecuaciones backward siempre son v\'alidas, mientras que las forward requieren condiciones de regularidad que se satisfacen en todos los ejemplos de este cap\'itulo.
\end{Observacion}

\section{Ejercicios}

\begin{Ejercicio}
  Supongamos que el n\'umero de llamadas por hora que llegan a un servicio de atenci\'on sigue un proceso de Poisson con $\lambda = 4$.
  \begin{enumerate}
    \item ¿Cu\'al es la probabilidad de que lleguen menos de dos llamadas en la primera hora?
    \item Supongamos que llegan seis llamadas en la primera hora. ¿Cu\'al es la probabilidad de que lleguen al menos dos llamadas en la segunda hora?
    \item La persona que atiende los tel\'efonos espera hasta que hayan llegado quince llamadas antes de ir a almorzar. ¿Cu\'al es el tiempo esperado que esperar\'a la persona?
    \item Supongamos que se sabe que exactamente ocho llamadas llegaron en las primeras dos horas. ¿Cu\'al es la probabilidad de que exactamente cinco de ellas llegaran en la primera hora?
    \item Supongamos que se sabe que exactamente $k$ llamadas llegaron en las primeras cuatro horas. ¿Cu\'al es la probabilidad de que exactamente $j$ de ellas llegaran en la primera hora?
  \end{enumerate}
\end{Ejercicio}

\begin{Ejercicio}
  Sean $X_t$ y $Y_t$ dos procesos de Poisson independientes con par\'ametros $\lambda_1$ y $\lambda_2$ respectivamente, que miden el n\'umero de clientes que llegan a las tiendas 1 y 2.
  \begin{enumerate}
    \item ¿Cu\'al es la probabilidad de que llegue un cliente a la tienda 1 antes de que llegue alg\'un cliente a la tienda 2?
    \item ¿Cu\'al es la probabilidad de que en la primera hora llegue un total de exactamente cuatro clientes a las dos tiendas?
    \item Dado que exactamente cuatro clientes han llegado a las dos tiendas, ¿cu\'al es la probabilidad de que los cuatro hayan ido a la tienda 1?
    \item Sea $T$ el tiempo de llegada del primer cliente a la tienda 2. Entonces $X_T$ es el n\'umero de clientes en la tienda 1 en el momento de la primera llegada a la tienda 2. Encuentre la distribuci\'on de probabilidad de $X_T$.
  \end{enumerate}
\end{Ejercicio}

\begin{Ejercicio}
  Supongamos que $X_t$ y $Y_t$ son procesos de Poisson independientes con par\'ametros $\lambda_1$ y $\lambda_2$, que miden el n\'umero de llamadas que llegan a dos tel\'efonos diferentes. Sea $Z_t = X_t + Y_t$.
  \begin{enumerate}
    \item Demuestre que $Z_t$ es un proceso de Poisson. ¿Cu\'al es su par\'ametro?
    \item ¿Cu\'al es la probabilidad de que la primera llamada llegue al primer tel\'efono?
    \item Sea $T$ el primer momento en que ha llegado al menos una llamada de cada uno de los dos tel\'efonos. Encuentre la funci\'on de densidad y la funci\'on de distribuci\'on de $T$.
  \end{enumerate}
\end{Ejercicio}

\begin{Ejercicio}
  Sea $\A$ el generador infinitesimal de una cadena de Markov irreducible en tiempo continuo con espacio de estados finito.
  \begin{enumerate}
    \item Sea $a$ un n\'umero positivo mayor que todas las entradas de $\A$. Demuestre que $P = (1/a)\A + \I$ es la matriz de transici\'on de una cadena en tiempo discreto irreducible y aperi\'odica.
    \item Use esto para concluir que $\A$ tiene un \'unico eigenvector izquierdo con eigenvalor 0 que es un vector de probabilidad, y que todos los dem\'as eigenvalores de $\A$ tienen parte real estrictamente menor que 0.
  \end{enumerate}
\end{Ejercicio}

\begin{Ejercicio}
  Sea $X_t$ una cadena de Markov con espacio de estados $\{1,2\}$ y tasas $\alpha(1,2) = 1$, $\alpha(2,1) = 4$. Encuentre $\PP_t$.
\end{Ejercicio}

\begin{Ejercicio}
  Repita el Ejercicio 3.5 con espacio de estados $\{1,2,3\}$ y tasas $\alpha(1,2) = 1$, $\alpha(2,1) = 4$, $\alpha(2,3) = 1$, $\alpha(3,2) = 4$, $\alpha(1,3) = 0$, $\alpha(3,1) = 0$.
\end{Ejercicio}

\begin{Ejercicio}
  Sea $X_t$ una cadena de Markov irreducible en tiempo continuo. Demuestre que para cada $i, j$ y todo $t > 0$,
  \[ \Prob\{X_t = j \mid X_0 = i\} > 0. \]
\end{Ejercicio}

\begin{Ejercicio}
  Considere la cadena de Markov en tiempo continuo con espacio de estados $\{1,2,3,4\}$ y generador infinitesimal
  \[ \A = \begin{pmatrix}
    -3 & 1 & 1 & 1 \\
    2 & -4 & 1 & 1 \\
    2 & 1 & -4 & 1 \\
    1 & 1 & 1 & -3
  \end{pmatrix}. \]
  \begin{enumerate}
    \item Encuentre la distribuci\'on de equilibrio $\bar{\pi}$.
    \item Suponga que la cadena comienza en el estado 1. ¿Cu\'al es el tiempo esperado hasta que cambie de estado por primera vez?
    \item Suponga nuevamente que la cadena comienza en el estado 1. ¿Cu\'al es el tiempo esperado hasta que la cadena est\'e en el estado 4?
  \end{enumerate}
\end{Ejercicio}

\begin{Ejercicio}
  Repita el Ejercicio 3.8 con
  \[ \A = \begin{pmatrix}
    -3 & 2 & 1 & 0 \\
    1 & -4 & 2 & 1 \\
    0 & 2 & -4 & 2 \\
    0 & 0 & 1 & -1
  \end{pmatrix}. \]
\end{Ejercicio}

\begin{Ejercicio}
  Suponga que $\alpha$ da las tasas para una cadena de Markov irreducible en tiempo continuo en un espacio de estados finito. Suponga que la medida invariante es $\pi$. Sea
  \[ p(x,y) = \alpha(x,y)/\alpha(x), \quad x \neq y, \]
  la probabilidad de transici\'on para la cadena en tiempo discreto correspondiente a la cadena en tiempo continuo "cuando se mueve". Encuentre la probabilidad invariante para la cadena en tiempo discreto en t\'erminos de $\pi$ y $\alpha$.
\end{Ejercicio}

\begin{Ejercicio}
  Sea $X_t$ un proceso de nacimiento y muerte en tiempo continuo con tasa de nacimiento $\lambda_n = 1 + 1/(n+1)$ y tasa de muerte $\mu_n = 1$. ¿Es este proceso positivamente recurrente, nulamente recurrente o transitorio? ¿Y si $\lambda_n = 1 - 1/(n+2)$?
\end{Ejercicio}

\begin{Ejercicio}
  Considere el modelo de poblaci\'on (Ejemplo 3, Secci\'on 3.3). ¿Para qu\'e valores de $\mu$ y $\lambda$ es segura la extinci\'on, es decir, cu\'ando es la probabilidad de alcanzar el estado 0 igual a 1?
\end{Ejercicio}

\begin{Ejercicio}
  Considere el modelo de poblaci\'on con inmigraci\'on (Ejemplo 4, Secci\'on 3.3). ¿Para qu\'e valores de $\mu, \lambda, \nu$ es la cadena positivamente recurrente, nulamente recurrente o transitoria?
\end{Ejercicio}

\begin{Ejercicio}
  Considere un proceso de nacimiento y muerte con $\lambda_n = 1/(n+1)$ y $\mu_n = 1$. Demuestre que el proceso es positivamente recurrente y encuentre la distribuci\'on estacionaria.
\end{Ejercicio}

\begin{Ejercicio}
  Suponga que se tiene un modelo determinista para la poblaci\'on donde la poblaci\'on crece proporcionalmente al cuadrado de la poblaci\'on actual. En otras palabras, la poblaci\'on $p(t)$ satisface la ecuaci\'on diferencial
  \[ \frac{dp}{dt} = c[p(t)]^2, \]
  para alguna constante $c > 0$. Suponga $p(0) = 1$. Resuelva esta ecuaci\'on diferencial (por separaci\'on de variables) y describa qu\'e sucede cuando el tiempo aumenta.
\end{Ejercicio}

\begin{Ejercicio}
  Considere un proceso de nacimiento y muerte con tasas de nacimiento $\lambda_n$ y tasas de muerte $\mu_n$. ¿Cu\'ales son las ecuaciones backward para las probabilidades de transici\'on $p_t(m,n)$?
\end{Ejercicio}

\end{document}